\subsection{Surfaces of general type} It is natural to wonder how far these techniques take us. We show

\begin{Theorem} \label{gtyp}Assume $X$ is a smooth projective minimal surface of general type. Let $L$ be a~$(k-1)$-very ample line bundle such that
\begin{gather*}
\chi(L)\geq 3k, \qquad L.K_X\geq 2K_X^2+k+1.
\end{gather*}
Then $L^{[k]}$ is big and nef over $X^{[k]}$.
\end{Theorem}

\begin{proof}
Arguing as in the proof of Theorem \ref{blok}, in particular equation \eqref{cvvv}, we express the Segre integral as the $t^k$-coefficient in the series
\begin{gather*}
f(t)=2^{-m} \left(\sqrt{1+2t}+\sqrt{1+6t}\right)^m {(1+2t)}^{\frac{n-1}{2}} (1+6t)^{\frac{p-1}{2}},
\end{gather*}
where
\begin{gather*}
m=L.K-2K^2, \qquad n=(L-K)^2+3\chi - 4k - 1, \qquad p=K^2-\chi+3,
\end{gather*}
with $K=K_X$ and $\chi=\chi(\mathcal O_X)$. Note that
\begin{gather*}
m+n+p=2\chi(L)-4k+2
\end{gather*}
is even. Furthermore, $p\geq 0$. Indeed, let $p_g$, $q$ denote the genus and irregularity of $X$. By~Noet\-her's inequality $K^2\geq 2p_g-4$, and $K^2>0$ since $X$ is minimal of general type. Averaging, we obtain $K^2>p_g-2$, and thus
\begin{gather*}
p=K^2-\chi+3>(p_g-2)-(1-q+p_g)+3=q\geq 0.
\end{gather*}
By Lemma \ref{cl}, the coefficients of
$f(t)$ up to order $\min\big(\frac{1}{2} (m+n+p)-1, m-1\big)$ are positive. In~particular, if
\begin{gather*}
k\leq \min\bigg(\frac{1}{2} (m+n+p)-1, m-1\bigg),
\end{gather*}
then the Segre integral is positive. The latter inequality is exactly our hypothesis, and thus $L^{[k]}$ is big and nef.
\end{proof}

